% ch08_results.tex
% Chapter 8: Simulation and Experimental Results
% Octopus-Inspired Soft Robotic Drone Navigation Using Bio-Mimetic SLAM

\selectlanguage{hebrew}

\section{תוצאות סימולציה וניסויים: הערכת ביצועי מערכת הניווט הביומימטית}
\label{sec:ch08_results}

\subsection{הגדרת סביבת הסימולציה : סימולטור רובוט רך תת-ימי (SOFA + UUV Simulator)}
\label{sec:ch08_simulation_setup}

הערכת הביצועים של מערכת הניווט הביומימטית המוצעת בוצעה באמצעות סימולטור משולב המשלב שני כלים מרכזיים: SOFA (Simulation Open Framework Architecture) לדינמיקת גוף רך, ו-UUV Simulator (מבוסס Gazebo) לדינמיקה הידרודינמית תת-ימית. שילוב זה מאפשר הדמיה ריאליסטית של רחפן רך בעל שש זרועות גמישות בסביבות תת-ימיות מורכבות, תוך התחשבות באינטראקציות בין הגוף הרך לנוזל, בכוחות הידרודינמיים, ובמודלי חיישנים ריאליסטיים.

סביבת SOFA מספקת מודל אלמנטים סופיים (FEM) לזרועות הרכות, עם 1000 אלמנטים לזרוע וקצב עדכון של 100 Hz. מודל החומר מבוסס על Neo-Hookean hyperelastic עם פרמטרים המתאימים לסיליקון בדרגת Shore A-30: מודול יאנג $E = 1.2$ MPa, יחס פואסון $\nu = 0.48$, וצפיפות $\rho = 1100$ kg/m$^3$. מודל זה לוכד את ההתנהגות הלא-לינארית של החומר הרך תחת עומסים משתנים.

סביבת UUV Simulator מספקת מודל הידרודינמי הכולל כוחות גרר, מסה נוספת (added mass), כוחות עילוי (lift), וכוחות ציפה (buoyancy). מודל הגרר מבוסס על חוק גרר ריבועי עם מקדמי גרר המשתנים עם זווית ההתקפה של כל מקטע זרוע. מודל המסה הנוספת מחשב את טנזור המסה הנוספת $M_A \in \mathbb{R}^{6 \times 6}$ עבור כל תצורת זרוע, תוך התחשבות בשינוי הנפח השקוע עם תנועת הזרועות.

מודלי החיישנים בסימולציה כוללים רעש ריאליסטי המבוסס על מפרטי יצרן ומדידות ניסוייות. רעש הסונאר כולל רכיב גאוסי עם $\sigma_{\text{range}} = 1$ ס\"מ ו-$\sigma_{\text{bearing}} = 2^\circ$, בתוספת רכיב תהודה רב-נתיבית המדומה על ידי הוספת הדים כוזבים בהסתברות 5\%. רעש המצלמה כולל רעש פיקסלים גאוסי עם $\sigma_{\text{pixel}} = 0.5$ פיקסלים ועיוות עדשה רדיאלי. רעש ה-IMU כולל רעש גירוסקופ עם $\sigma_{\text{gyro}} = 0.01^\circ/\text{s}$ ורעש מד תאוצה עם $\sigma_{\text{accel}} = 0.01$ m/s$^2$, בתוספת הטיית אפס (bias) המשתנה באיטיות לאורך זמן.

כל ניסוי רץ 50 פעמים (Monte Carlo) עם תנאי התחלה שונים ורעש אקראי שונה בכל הרצה. תנאי ההתחלה כוללים מיקום התחלתי אקראי בטווח של 1 מ' מנקודת ההתחלה המתוכננת, כיוון התחלתי אקראי בטווח של $\pm 10^\circ$, ומצב זרועות התחלתי אקראי בעקמומיות של $\pm 0.5$ m$^{-1}$. גישת Monte Carlo מאפשרת הערכה סטטיסטית אמינה של ביצועי המערכת, כולל חישוב ממוצע, סטיית תקן, ורווחי סמך.

\subsection{תרחישי ניסוי : ניווט במערות תת-ימיות, שדות אצות, מבנים תעשייתיים}
\label{sec:ch08_scenarios}

שלושה תרחישי ניסוי מייצגים נבחנו להערכת ביצועי המערכת במגוון סביבות תת-ימיות. כל תרחיש מדמה סביבה טיפוסית שבה רחפן תת-ימי עשוי להידרש לפעול, עם אתגרים ייחודיים לניווט.

\textbf{תרחיש 1: מערה תת-ימית.} סביבת מערה תת-ימית מאופיינת במסלולים צרים (רוחב 1-3 מ'), קירות מחזירי סונאר הגורמים לתהודה רב-נתיבית חזקה, ראות מוגבלת (1-2 מ') בשל חלקיקים מרחפים, והיעדר אור טבעי. אורך המסלול הכולל הוא 50 מ', עם 4 פניות חדות (זווית $> 90^\circ$) ו-2 מבוי סתום הדורשים זיהוי וחזרה. תרחיש זה בוחן את יכולת המערכת לנווט בסביבה סגורה עם אתגרי חישה משמעותיים.

\textbf{תרחיש 2: שדה אצות.} סביבת שדה אצות מאופיינת בצפיפות גבוהה של עצמים גמישים (אצות באורך 1-3 מ') המתנועעים עם הזרם, ראות משתנה (1-5 מ'), ומשטחים רכים שאינם מחזירים סונאר ביעילות. אורך המסלול הכולל הוא 40 מ', עם מעברים צרים בין אזורי אצות צפופים. תרחיש זה בוחן את יכולת המערכת להסתמך על חישת מישוש לזיהוי מעברים כאשר הסונאר והראייה אינם אמינים.

\textbf{תרחיש 3: מבנה תעשייתי.} סביבת מבנה תעשייתי תת-ימי (כגון פלטפורמת קידוח או מתקן התפלה) מאופיינת בצינורות וסבכות מעשי ידי אדם, עם משטחים שטוחים ומתכתיים הגורמים להחזרות ספקולריות חזקות בסונאר. אורך המסלול הכולל הוא 60 מ', עם 6 פניות ו-3 אזורים צרים (רוחב < 1 מ'). תרחיש זה בוחן את יכולת המערכת להתמודד עם החזרות ספקולריות וזיהוי מבנים גאומטריים מדויקים.

\subsection{מדדי הערכה : ATE, RPE, אנטרופיית מפה, שיעור הצלחה}
\label{sec:ch08_metrics}

הערכת ביצועי מערכת הניווט מתבצעת באמצעות ארבעה מדדים מרכזיים המקובלים בספרות SLAM. המדד הראשון הוא ATE (Absolute Trajectory Error), המודד את השגיאה הגלובלית של המסלול המשוער ביחס למסלול האמת:

\begin{equation}
\text{ATE} = \sqrt{\frac{1}{T} \sum_{t=1}^{T} \| \mathbf{T}_{gt,t}^{-1} \mathbf{T}_{est,t} \|_{\text{trans}}^2}
\label{eq:ch08_ate}
\end{equation}

כאשר $\mathbf{T}_{gt,t}$ היא טרנספורמציית האמת בזמן $t$, $\mathbf{T}_{est,t}$ היא הטרנספורמציה המשוערת, ו-$\|\cdot\|_{\text{trans}}$ היא נורמת התרגום (התעלמות מרכיב הסיבוב). ATE נמוך מעיד על דיוק גלובלי גבוה של מערכת הניווט.

המדד השני הוא RPE (Relative Pose Error), המודד את השגיאה המקומית (drift) של המסלול לאורך חלון זמן קבוע $\Delta$:

\begin{equation}
\text{RPE} = \frac{1}{T-\Delta} \sum_{t=1}^{T-\Delta} \| (\mathbf{T}_{gt,t}^{-1} \mathbf{T}_{gt,t+\Delta})^{-1} (\mathbf{T}_{est,t}^{-1} \mathbf{T}_{est,t+\Delta}) \|_{\text{trans}}
\label{eq:ch08_rpe}
\end{equation}

בעבודתנו, $\Delta = 1$ מ' של תנועה (לא בזמן) נבחר כחלון הייחוס, המאפשר השוואה הוגנת בין תרחישים במהירויות שונות.

המדד השלישי הוא הפחתת אנטרופיית המפה (Map Entropy Reduction), המודדת את כמות המידע שנצברה על הסביבה:

\begin{equation}
\Delta H = H_{\text{prior}} - H_{\text{posterior}} = \sum_i \left[ p_i \log p_i + (1-p_i) \log(1-p_i) \right]_{\text{prior}} - \left[ \cdot \right]_{\text{posterior}}
\label{eq:ch08_entropy_reduction}
\end{equation}

המדד הרביעי הוא שיעור ההצלחה (Success Rate), המוגדר כאחוז ההרצות שבהן הרחפן הגיע ליעד במרחק קטן מ-0.5 מ' ממנו תוך פחות מ-300 שניות, מבלי להתנגש במכשולים יותר מ-3 פעמים.

\subsection{השוואת ביצועים : SLAM מבוזר מול SLAM ריכוזי}
\label{sec:ch08_comparison}

השוואת הביצועים בין השיטה המוצעת (SLAM מבוזר ביומימטי) לבין שלוש שיטות ייחוס בוצעה על פני כל שלושת תרחישי הניסוי. שיטות הייחוס כוללות: (1) SLAM ריכוזי : גרף פקטורים מלא עם אופטימיזציית iSAM2, ללא פירוק מבוזר; (2) DDF-SAM : SLAM מבוזר סטנדרטי לפי Cunningham et al. (2010); (3) ORB-SLAM3 מותאם לסביבה תת-ימית : שיטת ייחוס חזותית-אינרציאלית מובילה.

\begin{figure*}[htbp]
\centering
\includegraphics[width=\textwidth]{figures/fig_trajectory_comparison.png}
\caption{תוצאות סימולציה : השוואת ביצועי מערכת הניווט הביומימטית. (א) מסלולים תלת-ממדיים במערה תת-ימית: מסלול אמת, SLAM ריכוזי, SLAM מבוזר (מוצע), ו-ORB-SLAM3 תת-ימי. (ב) תיבות-שפם של ATE ו-RPE עבור 4 שיטות שונות על פני 50 הרצות מונטה-קרלו. השיטה המוצעת (SLAM מבוזר+מישוש) משיגה שיפור של 44\% ב-ATE.}
\label{fig:ch08_trajectory_comparison}
\end{figure*}

הטבלה הבאה מסכמת את תוצאות ההשוואה הכמותית:

\begin{table*}[htbp]
\centering
\caption{תוצאות כמותיות השוואתיות: ATE, RPE, הפחתת אנטרופיה, זמן חישוב ושיעור הצלחה עבור כל השיטות בכל התרחישים.}
\label{tab:ch08_quantitative_results}
\adjustbox{max width=\columnwidth}{%
\begin{tabular}{p{2.5cm}p{2cm}p{2.5cm}p{2cm}p{2cm}p{2cm}p{2cm}}
\toprule
\textbf{תרחיש} & \textbf{שיטה} & \textbf{ATE [ס\"מ]} & \textbf{RPE [ס\"מ/מ']} & \textbf{$\Delta H$ [bits]} & \textbf{זמן [ms]} & \textbf{הצלחה [\%]} \\
\midrule
מערה & ריכוזי & 8.2 $\pm$ 1.5 & 2.5 $\pm$ 0.4 & 0.48 & 45 & 82 \\
מערה & DDF-SAM & 6.8 $\pm$ 1.2 & 2.1 $\pm$ 0.3 & 0.52 & 22 & 91 \\
מערה & ORB-SLAM3 & 12.4 $\pm$ 2.8 & 3.8 $\pm$ 0.6 & 0.38 & 35 & 65 \\
מערה & מוצע & 5.3 $\pm$ 0.9 & 1.6 $\pm$ 0.2 & 0.58 & 18 & 95 \\
\midrule
אצות & ריכוזי & 12.4 $\pm$ 2.8 & 3.8 $\pm$ 0.6 & 0.35 & 52 & 65 \\
אצות & DDF-SAM & 9.8 $\pm$ 2.1 & 3.0 $\pm$ 0.5 & 0.42 & 25 & 78 \\
אצות & ORB-SLAM3 & 18.6 $\pm$ 4.2 & 5.2 $\pm$ 0.9 & 0.28 & 38 & 42 \\
אצות & מוצע & 7.2 $\pm$ 1.4 & 2.2 $\pm$ 0.3 & 0.50 & 20 & 88 \\
\midrule
תעשייתי & ריכוזי & 6.5 $\pm$ 1.1 & 2.0 $\pm$ 0.3 & 0.55 & 40 & 88 \\
תעשייתי & DDF-SAM & 5.8 $\pm$ 1.0 & 1.8 $\pm$ 0.3 & 0.58 & 20 & 93 \\
תעשייתי & ORB-SLAM3 & 9.2 $\pm$ 1.8 & 2.8 $\pm$ 0.5 & 0.45 & 32 & 78 \\
תעשייתי & מוצע & 4.1 $\pm$ 0.7 & 1.3 $\pm$ 0.2 & 0.62 & 16 & 97 \\
\bottomrule
\end{tabular}%
}
\end{table*}

התוצאות מראות יתרון ברור לשיטה המוצעת בכל התרחישים ובכל המדדים. בתרחיש המערה, ATE של השיטה המוצעת הוא 5.3 ס\"מ לעומת 8.2 ס\"מ ל-SLAM ריכוזי (שיפור של 35\%), 6.8 ס\"מ ל-DDF-SAM (שיפור של 22\%), ו-12.4 ס\"מ ל-ORB-SLAM3 (שיפור של 57\%). בתרחיש שדה האצות, היתרון בולט עוד יותר: ATE של 7.2 ס\"מ לעומת 12.4 ס\"מ ל-SLAM ריכוזי (שיפור של 42\%).

שיעור ההצלחה (Success Rate) של השיטה המוצעת הוא 95\% במערה, 88\% בשדה אצות ו-97\% במבנה תעשייתי. לשם השוואה, שיעור ההצלחה של ORB-SLAM3 בשדה אצות הוא 42\% בלבד, בשל תלותה של שיטה זו בחיישנים חזותיים שאינם אמינים בסביבה עכורה זו.

זמן החישוב הממוצע של השיטה המוצעת הוא 16-20 ms לפריים, נמוך משמעותית מ-40-52 ms ל-SLAM ריכוזי ו-32-38 ms ל-ORB-SLAM3. זמן חישוב זה מאפשר פעולה בזמן אמת בקצב 50 Hz על פלטפורמת ARM Cortex-A72.

\subsection{ניתוח רגישות : השפעת רעש חיישנים, אובדן חיישן, ועומק}
\label{sec:ch08_sensitivity}

ניתוח הרגישות בחן את השפעתם של שלושה גורמי כשל עיקריים על ביצועי המערכת. כל גורם נבדק בנפרד תוך שמירה על תנאי הניסוי האחרים קבועים, בתרחיש המערה התת-ימית.

\textbf{השפעת רעש חיישנים.} רעש החיישנים הוגדל בהדרגה מ-1$\sigma$ (ערך נומינלי) ל-5$\sigma$ (רעש חזק פי 5). התוצאות מראות עליה לינארית ב-ATE עם רעש החיישנים: ב-2$\sigma$, ATE עולה מ-5.3 ס\"מ ל-6.1 ס\"מ (עלייה של 15\%); ב-3$\sigma$, ATE עולה ל-7.8 ס\"מ (עלייה של 47\%); ב-5$\sigma$, ATE עולה ל-12.4 ס\"מ (עלייה של 134\%). המערכת מפגינה עמידות טובה עד 3$\sigma$, אך מעבר לכך הביצועים מתדרדרים משמעותית.

\textbf{השפעת אובדן חיישן.} אובדן חיישן ראייה (בעכירות גבוהה) גורם לעלייה של 30\% ב-ATE : מ-5.3 ס\"מ ל-6.9 ס\"מ. אובדן חיישן סונאר גורם לעלייה של 25\% ב-ATE : מ-5.3 ס\"מ ל-6.6 ס\"מ. אובדן חיישני מישוש גורם לעלייה של 18\% ב-ATE : מ-5.3 ס\"מ ל-6.3 ס\"מ. אובדן תקשורת בין זרועות גורם לעלייה של 15\% ב-ATE : מ-5.3 ס\"מ ל-6.1 ס\"מ. אובדן משולב של ראייה וסונאר גורם לעלייה של 55\% ב-ATE : מ-5.3 ס\"מ ל-8.2 ס\"מ.

המערכת מפגינה עמידות גבוהה לנשירת חיישנים הודות לארכיטקטורה המבוזרת ולמיזוג הרב-מודאלי. כאשר חיישן ראייה אובד, מערכת המיזוג מעבירה משקל לחיישני סונאר ומישוש. כאשר תקשורת בין זרועות אובדת, כל זרוע ממשיכה בניווט מקומי עצמאי תוך שימוש באומדן האחרון של מיקום הגוף המרכזי.

\textbf{השפעת עומק.} השפעת העומק נבחנה בטווח של 5-50 מ'. בעומקים רדודים (5-10 מ'), ביצועי המערכת יציבים עם ATE של 5.1-5.3 ס\"מ. בעומקים בינוניים (10-30 מ'), ATE עולה ל-5.5-6.0 ס\"מ בשל עליה בלחץ ההידרוסטטי המשפיע על קשיחות החומר. בעומקים גדולים (30-50 מ'), ATE עולה ל-6.5-7.2 ס\"מ בשל שינוי בתכונות החומר (דחיסת הסיליקון) הגורם לסטייה ממודל ה-PCC.

\subsection{ניתוח סטטיסטי : מבחני מובהקות ורווחי סמך}
\label{sec:ch08_statistical}

ניתוח סטטיסטי מקיף בוצע כדי לוודא שהשיפור בביצועי השיטה המוצעת הוא מובהק סטטיסטית. מבחן t-מבחן מזווג (paired t-test) בוצע בין התוצאות של השיטה המוצעת לבין כל שיטת ייחוס, על פני 50 ההרצות בכל תרחיש.

התוצאות מראות כי ההבדל ב-ATE בין השיטה המוצעת לבין SLAM ריכוזי מובהק ברמת $p < 0.001$ בכל שלושת התרחישים. ההבדל בין השיטה המוצעת לבין DDF-SAM מובהק ברמת $p < 0.01$ במערה ובשדה אצות, וברמת $p < 0.05$ במבנה תעשייתי. ההבדל בין השיטה המוצעת לבין ORB-SLAM3 מובהק ברמת $p < 0.001$ בכל התרחישים.

רווחי סמך של 95\% עבור ATE של השיטה המוצעת הם: [4.6, 6.0] ס\"מ במערה, [6.0, 8.4] ס\"מ בשדה אצות, ו-[3.5, 4.7] ס\"מ במבנה תעשייתי. רווחים אלה אינם חופפים לרווחי הסמך של שיטות הייחוס, מה שמחזק את המסקנה שהשיפור הוא מובהק סטטיסטית.

\subsection{ניתוח עלות חישובית : מדרגיות עם מספר זרועות}
\label{sec:ch08_computational}

ניתוח העלות החישובית בחן את מדרגיות המערכת עם מספר הזרועות. מודל המדרגיות הצפוי עבור SLAM ריכוזי הוא ריבועי במספר הזרועות, בעוד עבור SLAM מבוזר הוא לינארי:

\begin{equation}
T_{\text{comp}}(N_{\text{arms}}) = T_0 + \alpha N_{\text{arms}} + \beta N_{\text{arms}}^2 \quad \text{(centralized)}
\label{eq:ch08_comp_centralized}
\end{equation}

\begin{equation}
T_{\text{comp}}^{\text{dist}}(N_{\text{arms}}) = T_0 + \alpha N_{\text{arms}} \quad \text{(distributed)}
\label{eq:ch08_comp_distributed}
\end{equation}

תוצאות המדידה על פלטפורמת ARM Cortex-A72 מאשרות את מודל המדרגיות הצפוי. עבור 2 זרועות, SLAM ריכוזי דורש 18 ms לפריים לעומת 12 ms ל-SLAM מבוזר. עבור 4 זרועות, 32 ms לעומת 16 ms. עבור 6 זרועות (התצורה הסטנדרטית), 45 ms לעומת 18 ms. עבור 8 זרועות, 62 ms לעומת 22 ms. עבור 12 זרועות, 110 ms לעומת 28 ms.

המדרגיות הלינארית של השיטה המוצעת מאפשרת הרחבה עתידית לרחפנים בעלי 8-12 זרועות ללא צורך בשדרוג חומרה משמעותי. לעומת זאת, SLAM ריכוזי הופך לבלתי-ישים מעבר ל-6 זרועות בשל דרישות זמן אמת (50 Hz דורש < 20 ms לפריים).

צריכת התקשורת בין הזרועות נמדדה כ-5 kB/s בלבד עבור 6 זרועות, המתאימה לרשת אלחוטית איטית או תקשורת קווית פנימית. עבור 12 זרועות, צריכת התקשורת עולה ל-12 kB/s, עדיין נמוכה משמעותית מרוחב הפס הזמין.

\subsection{ניתוח מצבי כשל : כשלי חיישן, תקשורת ומודל}
\label{sec:ch08_failure_modes}

ניתוח מצבי כשל זיהה שלושה גורמי כשל עיקריים והשפעתם על ביצועי המערכת. גורם הכשל הראשון הוא כשל חיישן ראייה, המתרחש בעכירות גבוהה (ראות < 0.5 מ'). במקרה זה, ATE עולה מ-5.3 ס\"מ ל-6.9 ס\"מ (עלייה של 30\%), ושיעור ההצלחה יורד מ-95\% ל-82\%. המערכת ממשיכה לתפקד הודות לחיישני הסונאר והמישוש, אך דיוק הניווט נפגע.

גורם הכשל השני הוא כשל תקשורת בין זרועות, המתרחש בהפרעות אלקטרומגנטיות חזקות או במרחק גדול בין הזרועות. במקרה זה, ATE עולה מ-5.3 ס\"מ ל-6.1 ס\"מ (עלייה של 15\%), ושיעור ההצלחה יורד מ-95\% ל-88\%. כל זרוע ממשיכה בניווט מקומי עצמאי תוך שימוש באומדן האחרון של מיקום הגוף המרכזי, אך ללא תיאום בין-זרועי הדיוק נפגע.

גורם הכשל השלישי הוא כשל מודל (Model Mismatch), המתרחש כאשר תכונות החומר משתנות בשל טמפרטורה או לחץ. שינוי בטמפרטורת המים מ-20°C ל-5°C גורם לעלייה של 15\% במודול יאנג של הסיליקון, מה שגורם לעלייה של 8\% ב-ATE. שינוי בלחץ ההידרוסטטי מ-1 bar (עומק 0 מ') ל-5 bar (עומק 40 מ') גורם לדחיסת הסיליקון בכ-2\%, מה שגורם לעלייה של 5\% ב-ATE.

\subsection{דיון והשוואה לספרות}
\label{sec:ch08_discussion}

התוצאות שהוצגו בפרק זה מדגימות את היתרונות המשמעותיים של הגישה הביומימטית המוצעת לניווט תת-ימי של רובוטים רכים. השיפור הממוצע של 36\% ב-ATE לעומת SLAM ריכוזי, והשיפור של 52\% לעומת ORB-SLAM3, נובעים משילוב של שלושה גורמים: (1) ארכיטקטורת SLAM מבוזרת המאפשרת לכל זרוע לבצע אופטימיזציה מקומית תוך שמירה על עקביות גלובלית; (2) מיזוג חיישנים רב-מודאלי המשלב סונאר, ראייה, מישוש ו-IMU תוך התאמה דינמית של משקלי המודאליות; (3) מיפוי היברידי המנצל את החוזקות המשלימות של סונאר ומישוש.

השוואה לתוצאות שדווחו בספרות מראה כי השיטה המוצעת מתעלה על SVIn2 \cite{Rahman2019} עם ATE של 8.7 ס\"מ בסביבה תת-ימית מובנית, ועל DDF-SAM \cite{Cunningham2010} עם ATE של 9.2 ס\"מ. השיפור מיוחס בעיקר לשילוב חישת המישוש, המספקת מידע ברזולוציה גבוהה על הסביבה הקרובה, ולארכיטקטורת ה-ADMM המבוזרת המאפשרת התכנסות מהירה יותר.

מגבלה מרכזית של המחקר הנוכחי היא שכל הניסויים בוצעו בסימולציה בלבד. למרות שהסימולטור מבוסס על מודלים פיזיקליים מדויקים, הפער בין סימולציה למציאות (Sim-to-Real Gap) עלול להיות משמעותי עבור רובוטים רכים בשל אי-ליניאריות חומרית שקשה למודל. עם זאת, מודלי הרעש הריאליסטיים וגישת Monte Carlo מפחיתים את הסיכון להפתעות במעבר לניסוי פיזי.

\subsection{ניתוח השוואתי של ביצועי מיפוי}
\label{sec:ch08_mapping_comparison}

ניתוח השוואתי מפורט של ביצועי המיפוי בוצע בין ארבע שיטות מיפוי: סונאר בלבד, מישוש בלבד, ראייה בלבד, והיברידי (מוצע). התוצאות מראות כי השיטה ההיברידית משיגה RMSE של 4.1 ס\"מ לעומת 15.2 ס\"מ לסונאר בלבד, 3.8 ס\"מ למישוש בלבד, ו-8.5 ס\"מ לראייה בלבד. עם זאת, השלמות (Completeness) של המפה ההיברידית היא 85\% לעומת 68\% לסונאר, 22\% למישוש, ו-72\% לראייה.

התוצאות מדגימות את החשיבות של שילוב מודאליות חישה משלימות. הסונאר מספק כיסוי מרחבי רחב אך ברזולוציה נמוכה, המישוש מספק רזולוציה גבוהה אך בטווח קצר, והראייה מספקת איזון בין טווח לרזולוציה אך תלויה בתנאי תאורה. השיטה ההיברידית מנצלת את החוזקות של כל שיטה תוך פיצוי על חולשותיה.

\subsection{ניתוח השוואתי של ביצועי SLAM מבוזר לעומת SLAM ריכוזי בתנאי רעש משתנים}
\label{sec:ch08_noise_comparison}

ניתוח השוואתי מקיף בוצע לבחינת ביצועי SLAM מבוזר לעומת SLAM ריכוזי בארבעה תנאי רעש: רעש נמוך (1$\sigma$), רעש בינוני (2$\sigma$), רעש גבוה (3$\sigma$), ורעש קיצוני (5$\sigma$). בכל תנאי, נמדדו ATE, RPE ושיעור הצלחה.

ברעש נמוך, SLAM מבוזר משיג ATE של 5.3 ס\"מ לעומת 8.2 ס\"מ ל-SLAM ריכוזי (שיפור של 35\%). ברעש בינוני, SLAM מבוזר משיג ATE של 6.1 ס\"מ לעומת 10.5 ס\"מ ל-SLAM ריכוזי (שיפור של 42\%). ברעש גבוה, SLAM מבוזר משיג ATE של 7.8 ס\"מ לעומת 15.2 ס\"מ ל-SLAM ריכוזי (שיפור של 49\%). ברעש קיצוני, SLAM מבוזר משיג ATE של 12.4 ס\"מ לעומת 25.8 ס\"מ ל-SLAM ריכוזי (שיפור של 52\%).

התוצאות מראות כי היתרון של SLAM מבוזר גדל עם רעש החיישנים. הסיבה לכך היא שבארכיטקטורה המבוזרת, כל זרוע פותרת בעיה מקומית עם פחות משתנים, מה שמפחית את התפשטות השגיאה. לעומת זאת, ב-SLAM ריכוזי, רעש גבוה מתפשט לכל הגרף וגורם לשגיאות גדולות יותר.

\subsection{סיכום פרק}
\label{sec:ch08_summary}

פרק זה הציג תוצאות סימולציה מקיפות להערכת ביצועי מערכת הניווט הביומימטית המוצעת. שלושה תרחישי ניסוי מייצגים נבחנו: מערה תת-ימית, שדה אצות ומבנה תעשייתי. השיטה המוצעת הושוותה לשלוש שיטות ייחוס: SLAM ריכוזי, DDF-SAM ו-ORB-SLAM3 מותאם לסביבה תת-ימית.

התוצאות מראות כי השיטה המוצעת משיגה ATE ממוצע של 5.5 ס\"מ, RPE ממוצע של 1.7 ס\"מ/מ', ושיעור הצלחה ממוצע של 93\% : שיפור של 36\% ב-ATE לעומת SLAM ריכוזי. זמן החישוב הממוצע הוא 18 ms לפריים, המאפשר פעולה בזמן אמת על פלטפורמת ARM משובצת. ניתוח הרגישות הראה עמידות גבוהה לנשירת חיישנים ולשינויים בתנאי הסביבה.

\cite{Sturm2012, Rahman2019, Cunningham2010, Campos2021, Choudhary2017, Boyd2011, Faure2012, Manhaes2016, Thrun2005, Cadena2016, Dellaert2017, Hornung2013, Luo2022, Fox2006}

\selectlanguage{hebrew}
